\newpage
\phantomsection
\label{prf:two-sided-successor-predecessor-inverses}
\LRAProofFor{thm:two-sided-successor-predecessor-inverses}

\begin{remark*}[Return]
\hyperref[thm:two-sided-successor-predecessor-inverses]{Return to Theorem}
\end{remark*}

\begin{theorem*}[Successor and predecessor are inverse operations]
In the two-sided construction, successor and predecessor undo each other
everywhere, and both maps are injective.
\end{theorem*}

\begin{proof}[Professional Standard Proof]
\LRAProofBodyStart
The defining clauses for successor and predecessor are checked by cases on
the three parts of the carrier \(Z=\{0\}\sqcup P\sqcup N\). At \(0\), successor
enters the first positive point and predecessor returns to \(0\); predecessor
enters the first negative point and successor returns to \(0\). On the positive
ray, predecessor removes one successor step unless the point is the first
positive point, in which case it returns to \(0\). On the negative ray,
successor removes one predecessor step unless the point is the first negative
point, in which case it returns to \(0\). These are exactly the two identities.

If \(\mathsf{succ}(x)=\mathsf{succ}(y)\), applying
\(\mathsf{pred}\) to both sides gives \(x=y\). The proof for predecessor
is the same, applying successor to both sides.
\end{proof}

\begin{proof}[Detailed Learning Proof]
\LRAProofBodyStart
The point of the definition is that the only delicate places are the two
crossings through zero. Everywhere else, successor and predecessor simply add
or remove one constructor on a ray. Thus the proof is a structural case split:
zero, positive point, negative point. The injectivity statements are not new
axioms; they are consequences of having inverse functions.
\end{proof}

\begin{remark*}[Proof structure]
Case split on the carrier, then derive injectivity by composing with the
opposite operation.
\end{remark*}

\begin{dependencies}
\begin{itemize}
  \item \hyperref[def:two-sided-successor-construction-data]{Two-sided successor construction data}
\end{itemize}
\end{dependencies}

\clearpage
